\section{The board, the verifiers, and the human audit}
\label{sec:data}

Everything in this paper uses one fixed set of model responses that the leaderboard published.
This section describes those responses, the three verifiers that read them, how we sort the
responses by what the verifiers say, and the small human audit that tells us what a careful
reader would say.

\subsection{The board and the roster}

The Open LLM Leaderboard \citep{openllmleaderboard, fourrier2024openllm} tested submitted models
on MATH-Hard, the \NumItems{} hardest (Level~5) problems of the MATH benchmark of competition
mathematics \citep{hendrycksmath2021}, spread over seven subjects, and it published every model's
written response to every problem. The leaderboard has since been retired, which suits this
study: its responses can no longer change, and its scoring is on record. We call its table of
models the \emph{board}.

Our \emph{roster} is the \NumModels{} models we study, a small part of a board that held
\BoardModelsRescored{} models at the February 2025 re-scoring \citep{kydlicek2025fixing}. The
roster was not chosen for this paper. An earlier study ranked all \RosterCatalogued{} of the
leaderboard's per-model result repositories by download count and kept the top \RosterSelected{},
at most \RosterMaxPerOrg{} per organisation, under a rule written down before any re-scoring; the
roster is the \NumModels{} of these that received at least one problem in the human audit described
below (Appendix~\ref{app:data}). It leans towards popular models, is neither complete nor random,
and sits low on the board: the median roster model is at the \RosterPctMedian{}th percentile of
displayed accuracy (Appendix~\ref{app:board}). Every roster model answered every problem, so we
work with all \NumRows{} responses.

\subsection{Three verifiers}
\label{sec:verifiers}

A verifier gives each response a verdict of $1$ (credit) or $0$ (no credit) in two steps.
\emph{Extraction} looks for the answer in the text, and the \emph{equivalence check} then decides
whether that answer matches the gold answer. Either step can fail, and a response whose answer
cannot be found gets $0$, exactly like a wrong answer. Our three verifiers differ mainly in how they
extract.

\emph{The 2024 verifier} is the leaderboard's original scorer, from the lm-evaluation-harness
\citep{biderman2024lessons}. It reads an answer only from the sentence ``Final Answer: The final
answer is X. I hope it is correct.'', which the prompt demonstrates and which we call the
\emph{template}, and it checks whether that answer is mathematically equal to the gold answer
rather than identical as text. The leaderboard published this verifier's verdict on every response, together
with the scoring code that ran. Re-running that code reproduces the published verdict on
\EmReexecItems{} of \EmReexecItems{} responses, so we know exactly which verifier produced the
original scores.

\emph{Math-Verify} is the verifier the leaderboard switched to in February 2025
\citep{kydlicek2025fixing}. The switch changed the displayed scores but left the published
per-response verdicts as they were, so from then on the board showed scores from one verifier and
published verdicts from another. Running Math-Verify ourselves reproduces the re-scored totals in
\MvAggregateMatch{} of \NumSubjectsRescored{} model-subject combinations; each of the others is off
by a single problem.

\emph{The boxed verifier} reads the expression inside the last \texttt{\textbackslash boxed\{\}}
command, which is how the benchmark's own reference solutions mark their answers, and compares it
in the same way. Appendix Table~\ref{tab:provenance} gives each verifier's version and the check
behind it.

The three are not independent. All rest on the same algebra library, and the 2024 and boxed
verifiers also share the harness's normalisation of answers and the stored gold answer, which the
harness takes from the same reference solution that Math-Verify reads (Appendix
Table~\ref{tab:shared}). Nothing below assumes independence; what matters is that the verifiers fail
on different responses, and they do. While tracing the verdicts we also found that our own boxed
verifier had been silently degraded on Windows, where an error handler turned every check of
mathematical equality into a plain comparison of text. All results use the corrected version. Appendix~\ref{app:data}
reports the defect, which is itself an instance of this paper's subject: a verifier that fails
quietly still produces a plausible number.

\subsection{Sorting responses by the verdicts}

Writing the $K = \NumVerifiers$ verdicts in a fixed order (the 2024 verifier, the boxed verifier,
Math-Verify) gives each response a three-digit string, which places it in one of $\NumCells$
\emph{cells}.
Cell $111$ holds the responses that all three verifiers credit and cell $000$ those that none
credits; in cell $101$, for example, the 2024 verifier and Math-Verify credit the response and the
boxed verifier does not. We call $111$ and $000$ the \emph{concordant} cells, because there the
verifiers agree, and the other six the \emph{discordant} cells. Table~\ref{tab:cells} gives their
sizes.

\subsection{The human audit and the reader verdict}

The quantity this paper is after is whether the answer a response gives is right, as a careful
reader would judge it. A ranking can also legitimately reward following the requested format; we
measure that separately in Section~\ref{sec:disagree}. We obtain the reader's judgement in two
steps, and only the first is human. A labeller sees the problem and the response and nothing
else: not the model, not any verifier's verdict, and not the gold answer. The labeller copies down,
exactly as written, the final answer the response \emph{commits to}, meaning the answer it presents
as final after any corrections it makes along the way. If there is no such answer the labeller
writes \textsc{none}, and if the response gives two different final answers without choosing
between them, \textsc{contradictory}. The benchmark's own equivalence check then compares the copied
answer with the gold answer. The result is the \emph{reader verdict}, written $Y^*$: $1$ if they
match, and $0$ otherwise, including for \textsc{none} and \textsc{contradictory} (Appendix~\ref{sec:sensitivity}
drops those instead). $Y^*$ is a measurement, not mathematical truth: it records which answer a
reader finds and whether the benchmark accepts it.

All human labels come from one audit of 400 responses, designed for an earlier study in a dated
document committed before labelling began. The labellers worked from a spreadsheet that showed only
the problem and the response. Of the 400 labels, \AuditNAgreeAll{} fall in cell $111$ and
\AuditNRejectAll{} in cell $000$. Because the audit was drawn by the boxed verifier's verdict and by
subject, not by model and cell, it has too few labels to say anything about one model in one
discordant cell. We therefore use only its labels in the two concordant cells and treat every
discordant cell as unknown, which can only make our conclusions more cautious. The author labelled
all 400 responses; two other labellers labelled overlapping sets of \LabTwoN{} and \LabThreeN{},
and using their labels instead moves the certified count by at most \LabSwapMax{}
(Appendix~\ref{app:data}). Our analysis plan also describes a larger audit drawn by model and cell
and a complete labelling of the discordant responses; neither has been carried out, and no result
here depends on them.
